% ── Section 7: Conclusion ────────────────────────────────────────────
\section{Conclusion}
\label{sec:conclusion}

This survey examined prompt injection and tool-use attacks in LLM-integrated applications through a unified lens. We organized the threat space into three attack families---direct injection, indirect injection, and tool-use exploitation---and mapped four defense categories---input sanitization, instruction hierarchy, output monitoring, and architectural isolation. Our cross-cutting comparison revealed that existing defenses vary widely in their threat-model assumptions, evidence strength, and deployment readiness, with the Instruction Hierarchy being the only approach confirmed in production use. Two findings deserve particular emphasis: the positive correlation between model capability and injection susceptibility suggests that the problem will worsen as models improve, and the expansion of tool-use in agentic applications means that successful injections can now cause real-world harm, not just misleading text. We urge the research community to converge on standardized evaluation benchmarks, to explore architectural solutions that enforce data-instruction separation at the model level, and to study the emerging risk of injection propagation in multi-agent systems.
